% FULL_SOLUTION_SCOPE: {"target":"The entire universal four-total-insertion/deletion connectivity statement of catalogue ID 3100089, with the original pressing rule and pairwise connectivity convention.","result":"Theorem 1: vertices 1..5 form a cycle, vertex 6 is adjacent only to 2, and initially exactly 4 and 5 are black. Its complete successful-word family is {432651,512643}, with insertion-deletion distance eight; H4 is disconnected.","coverage":"One admissible finite simple colored graph with a nonempty successful-word family disproves the universal assertion. The proof exhausts every legal branch for this graph and excludes every possible intermediate successful word. Other graph sizes, colorings, empty families and the empty graph need not be classified to disprove a universal statement.","dependencies":"No external theorem or computational certificate is needed for the self-contained proof; only the original pressing and insertion/deletion definitions. Literature and graph-atlas data supplied context and discovery.","human_verification":"None. Wentao Li provided direction and prompts; no independent human proof checking occurred.","unverified":"The complete proof has not been independently reviewed by a human or formalized in a theorem prover. Repository checks and same-model audits do not supply that verification.","novelty":"Expanded primary-source and duplicate search on 2026-10-11 located no matching colored five-cycle-with-pendant example in inspected material. Bixby-Flint-Miklos explicitly conjecture eight total edits, and this result does not refute that formulation. Whitlatch dissertation full text and JMM full proof remain inaccessible; historical novelty is unestablished."}
\documentclass[11pt]{article}
\usepackage[T1]{fontenc}
\usepackage[margin=25mm]{geometry}
\usepackage{amsmath,amssymb,amsthm,booktabs,xurl,hyperref}
\newtheorem{theorem}{Theorem}
\newtheorem{corollary}[theorem]{Corollary}
\newcommand{\ed}{\operatorname{ed}}
\newcommand{\LCS}{\operatorname{LCS}}
\title{A six-vertex counterexample to\newline four-total-edit pressing connectivity}
\author{OpenAI GPT 6.1 Sol\\\small Max reasoning}
\date{11 October 2026}
\begin{document}
\maketitle
\begin{abstract}
The four-total-insertion/deletion connectivity assertion recorded as
UnsolvedMath ID 3100089 is false. A five-cycle with one pendant vertex,
colored black on two adjacent cycle vertices, has exactly two successful
pressing words. Their longest common subsequence has length two, although
both words have length six; their total insertion--deletion distance is
therefore eight. We give the entire branching proof directly from the
pressing rule. This is a full disproof of the recorded four-total-edit
assertion, while the separate formulation allowing four deletions from
\emph{each} word is consistent with this example. Supporting exhaustive
checks reproduce the proof. Historical novelty and independent human or
formal proof verification are not established.
\end{abstract}

\section{Original target and conventions}
The target is the unchanged repository definition of UnsolvedMath
ID 3100089 (AMR-030-0089), whose source-catalogue status remains
\emph{open} \cite{catalogue}. A bicolored graph is a finite simple graph
with each vertex black or white. A legal press at a currently black vertex
$v$ complements the edges between distinct vertices of its closed
neighborhood $N[v]$ and flips every color there. All other edges and
colors remain unchanged. A successful pressing word consists of legal
presses and ends with an edgeless all-white graph. In particular, an
all-white graph with a surviving edge is \emph{not} a successful terminal
state. A pressed vertex becomes isolated and white, and stays that way,
so no vertex can appear twice in a legal word.

Each insertion or deletion of one vertex symbol costs one edit. Thus a
replacement costs two; there is no unit-cost substitution operation.
Intermediate words in computing edit distance need not be successful.
Words distinguish the individual vertex symbols, including vertices
exchanged by a symmetry of the graph.
Let $H_4(G)$ have the successful words as vertices, with distinct words
joined when their total insertion--deletion distance is at most four.
The exact repository target, with the four-edit wording read literally,
is whether, for \emph{every} such graph and coloring, every pair of successful
words can be joined by a finite path in $H_4(G)$
\cite[Conjecture 1]{cooper}. The pairwise assertion is vacuous
for an empty successful-word family; it also holds for a singleton.
The empty word succeeds precisely for an initially edgeless all-white
graph. These conventions are preserved.

\section{Counterexample and complete proof}
Let $G$ have vertices $V=\{1,2,3,4,5,6\}$, edges
\[
 E=\bigl\{\{1,2\},\{2,3\},\{3,4\},\{4,5\},\{5,1\},\{2,6\}\bigr\},
 \qquad B=\{4,5\},
\]
where $B$ is the set of initially black vertices. Thus $1,2,3,4,5$ form a
five-cycle and $6$ is a white pendant vertex at $2$. This is a connected
finite simple graph and has black vertices, so the counterexample does
not rely on vacuous connectivity or on lack of successful pressing words.

\begin{center}
\begin{tabular}{lll}
\toprule
Vertex & Neighbors & Initial color\\
\midrule
$1$ & $2,5$ & white\\
$2$ & $1,3,6$ & white\\
$3$ & $2,4$ & white\\
$4$ & $3,5$ & black\\
$5$ & $1,4$ & black\\
$6$ & $2$ & white\\
\bottomrule
\end{tabular}
\end{center}

\begin{theorem}\label{counterexample}
The complete successful-word family of $G$ is
\[
 \sigma=4\,3\,2\,6\,5\,1,
 \qquad
 \tau=5\,1\,2\,6\,4\,3.
\]
Their total insertion--deletion distance is eight. Consequently $H_4(G)$
consists of two isolated vertices, and the original universal connectivity
assertion is false.
\end{theorem}
\begin{proof}
We first determine every successful word directly from the graph rule.
Write $ij$ for the unordered edge $\{i,j\}$, and write a pressing prefix
without separators since every symbol is a single digit. The color-preserving
automorphism
\[
             \rho=(1\ 3)(4\ 5),\qquad \rho(2)=2,\quad\rho(6)=6
\]
exchanges the two initially black vertices. Every legal nonempty word
starts with $4$ or $5$. It therefore suffices to analyze the branch
starting with $4$, and then apply $\rho$.

The complete half-tree is the following table. All six vertices remain
present; vertices already pressed are isolated and white. The right
column lists \emph{all} available legal next presses.
\begin{center}
\begin{tabular}{lll}
\toprule
Prefix & Remaining edges & Black vertices\\
\midrule
$\varnothing$ & $12,15,23,26,34,45$ & $4,5$\\
$4$ & $12,15,23,26,35$ & $3$\\
$43$ & $12,15,25,26$ & $2,5$\\
$435$ & $26$ & $1$\\
$4351$ & $26$ & none\\
$432$ & $16,56$ & $1,6$\\
$4321$ & $56$ & none\\
$4326$ & $15$ & $5$\\
$43265$ & none & $1$\\
$432651$ & none & none\\
\bottomrule
\end{tabular}
\end{center}

Here is also an explicit derivation of the branching rows. Pressing $4$
uses $N[4]=\{3,4,5\}$: the edges $34,45$ disappear, $35$ appears,
and only $3$ is black. Pressing that forced $3$ uses
$\{2,3,5\}$: $23,35$ disappear, $25$ appears, and the black set becomes
$\{2,5\}$. The remaining edges are $12,15,25,26$.

If $5$ is chosen next, its closed neighborhood is $\{1,2,5\}$.
All three edges of that triangle disappear, leaving only $26$ and the
black vertex $1$. Pressing the isolated $1$ leaves an all-white graph
with edge $26$, so this branch cannot succeed.

If $2$ is chosen after $43$, its closed neighborhood is
$\{1,2,5,6\}$. All edges incident with $2$ disappear; $15$ disappears
and $16,56$ appear. The black set becomes $\{1,6\}$. Pressing $1$
at this point removes $16$ and makes $6$ white, leaving the all-white
edge $56$; this branch also fails. Choosing $6$ instead uses
$\{1,5,6\}$, removes $16,56$, creates $15$, and leaves $5$ as the
only black vertex. The next press $5$ removes $15$ and makes $1$
black; pressing $1$ then gives the edgeless all-white graph.

All decisions have been exhausted: the only successful word starting
with $4$ is $\sigma$. Applying $\rho$ gives the only successful word
starting with $5$, namely $\tau$. The initial graph is not terminal,
so the empty word adds no success. A stalled all-white graph admits
no further legal press, and a successful terminal state also admits
no further press. This proves completeness of the two-word family,
without any matrix theorem or computational assumption.

For the distance, divide the symbols into the ordered pairs
\[
                 P=(4,3),\qquad Q=(2,6),\qquad R=(5,1).
\]
The words have block orders $\sigma=PQR$ and $\tau=RQP$, with the
same internal order within each pair. The relative order of any two
symbols from different blocks is reversed between the two words.
Thus a common subsequence can contain symbols from only one block,
and has length at most two. Each pair itself is common, so
$\LCS(\sigma,\tau)=2$.

Under insertions and deletions, the original symbols that survive
form a common subsequence. At most two of the six can survive, so
at least four deletions and four insertions are necessary. Conversely,
keep the common pair $26$, delete the other four symbols from
$\sigma$, and insert $51$ before it and $43$ after it. This gives
$\tau$ with eight edits. Hence $\ed(\sigma,\tau)=8$ exactly.
Since these are the only successful words, they have no edge in
$H_4(G)$ and there is no intermediate successful vertex through which
to connect them. This disproves the universal statement.
\end{proof}

\begin{corollary}\label{lowerbound}
Any total insertion--deletion threshold that connects successful words
for every bicolored graph must be at least eight. For the displayed graph,
eight is exactly the smallest threshold giving connectivity.
\end{corollary}
\begin{proof}
The same two-vertex word graph has no edge at any threshold below eight,
and has its single connecting edge at threshold eight.
\end{proof}

\section{Formulation and literature check}
The distinction between \emph{four total edits} and \emph{four deletions
from each word} is essential. Cooper--Davis \cite[Conjecture 1]{cooper}
and the repository definition use the former wording. Their 2015 revision
also expresses doubts about that statement for general graphs. The
counterexample above answers that precise four-total-edit question
negatively. It does not replace the original target with a different bound.

Bixby--Flint--Mikl\'os \cite[Conjecture 6]{bixby} use a longest-common-
subsequence deficit of four for the general conjecture. For two successful
words of common length $k$, a deficit $d$ means total insertion--deletion
distance $2d$. Their general formulation therefore permits eight total
edits. Their Theorem 7 proves the stronger deficit-two result for paths.
Our graph is not a path, and the two words here have deficit four,
so the example is consistent with both of those statements. These
literature results supply context and are not dependencies of our proof.

An indexed Joint Mathematics Meetings abstract by Bowen Li, Istv\'an
Mikl\'os, and Carter Rockhold Teplica \cite{jmm}, for January 2024,
announces connectivity when neighboring words can be made identical
by deleting at most four terms from \emph{each}. That is likewise
the eight-total-edit formulation. The indexed abstract was inspected;
direct access to the PDF failed and no underlying proof was accessed.
We make no claim to have verified that general announcement, and our
counterexample does not contradict it.

The Bixby--Flint--Mikl\'os arXiv formulation already states the
larger bound explicitly. This manuscript gives a counterexample to
the \emph{literal four-total-edit formulation} and clarifies
the difference between the recorded formulations. We do not claim to have
refuted Bixby--Flint--Mikl\'os Conjecture 6 or resolved the historical
Pressing Game Conjecture under its eight-total-edit meaning.

The pre-submission review on 11 October 2026 expanded the source check.
Besides the Cooper--Davis and Bixby--Flint--Mikl\'os papers, we inspected
the accessible text of the uniquely pressable graph and autonomous poset
papers \cite{unique,vposets}. The illustrated two-poset example in
\cite[Figure 2]{vposets} has four vertices, unlike our graph; its illustrated
four-vertex generators and the text's two-word example are also distinct.
We inspected the bicolored graph illustration in the Cooper--Whitlatch
poster \cite{poster}; it is a different six-vertex graph. Mikl\'os's
dissertation \cite[Chapter 9, pp.~167--178]{miklosthesis}, in its
conjecture statement and discussion, uses the LCS-deficit-four formulation
and the stronger path result. No matching counterexample was located in
these inspected passages or figures.

We also consulted Cooper's problems page \cite{problems} and the indexed
JMM abstract. Searches included counterexamples, two successful words,
six-vertex graphs, cycles with pendant vertices, equivalent edit/LCS
wordings, and the named conjecture and authors. Queries of the target
repository's public issues and pull requests for the ID and pressing
terminology returned no matching records at the time of checking.
These are search results, not a proof of absence of prior work.

Whitlatch's 2019 dissertation \cite{whitlatchthesis} is particularly
relevant because its abstract names the complexity of the pressing space.
Its landing page and abstract were accessible, but the linked full-text
PDF returned HTTP 403 and was not inspected. The JMM full proof was also
not accessed. Historical novelty therefore remains \emph{unestablished}.
Neither absence of a search result nor the size of the example justifies
calling it new. Prior publication of this exact graph, a relabeling, or an
equivalent construction is not ruled out.

\section{Discovery, reproducibility, and verification}
This new version preserves the earlier star manuscript and the earlier
double-star manuscript \cite{priorstar,priordouble} without alteration.
Those restricted positive results do not address this cyclic graph. They
motivated continued research, but no lemma from them is required for the
counterexample proof.

Discovery used a standard-library Python enumeration of successful
pressing words, with augmented-adjacency updates and common two-deletion
signatures for four-edit adjacency. The input was the NetworkX graph-atlas
data \cite{atlas}, downloaded from its public repository without importing
the NetworkX package. The search stopped at atlas index 104, the graph
above, with zero-based black-color mask 24. It had checked 4,596
graph/color cases and enumerated 97,571 successful words; none of those
cases was skipped for the word-count cap. This is a discovery log,
not an exhaustive check of the entire atlas or of all graphs. No
vertex-minimality claim is needed or made here.

The separate repository script
\begin{quote}\ttfamily\small
scripts/verify\_pressing\_3100089\_counterexample.py
\end{quote}
does not import the search script or require atlas data. It enumerates
the whole legal tree by edge sets and color sets: 19 prefixes, exactly
the two displayed successes, and four stalled maximal words
\[
              4321,\quad4351,\quad5123,\quad5143.
\]
It checks every row of the proof table and the reflection $\rho$.
A separate dense principal-determinant implementation checks all 64
vertex subsets and all 1,957 distinct-symbol words of lengths zero
through six, including all 720 full permutations. Its legal flags
agree exactly with the direct pressing tree, and its full successful
family agrees with the theorem. Direct insertion/deletion dynamic
programming and enumeration of common subsequences both confirm
distance eight and longest-common-subsequence length two. All these
checks passed with Python 3.13.2.

Run from the repository root, with assertions enabled:
\begin{verbatim}
python3 scripts/verify_pressing_3100089_counterexample.py \
  --output /absolute/path/to/counterexample-results.json
\end{verbatim}
The JSON retains the initial graph, complete state tree, terminal words,
counts, distance, software version, and verifier-source SHA-256.
The proof itself is self-contained; the script supports its arithmetic
and branch completeness, rather than being an essential finite certificate.

The mandatory initial Prompt A was recorded before this continuation's
research. Prompt B was applied to the complete proposed manuscript and
the corrected complete output was retained, together with the full audit,
under \texttt{research/pressing-3100089-continuation-2026-10-11/} in the
shared workspace, outside the contribution repository. The atlas data,
search source, candidate record, raw checks, and earlier manuscripts are
also retained. They have not yet been publicly archived. Public web
browsing, Python, Git, the repository structural validator, and the app
LaTeX compiler were used. Structural validation and compilation check
formatting, not mathematical correctness. No delegated agent, independent
human checker, external peer review, or theorem prover was used.

\paragraph{Pre-submission recheck and version history.}
Version 3 preserves version 2 locally and leaves its graph, Theorem 1, branching
table, and distance proof unchanged. The entire proof was checked again
against the original definition, including terminal success, every legal
choice, persistence of pressed vertices, the reflection, survivor ordering,
and the absence of intermediate successful words. The existing verifier
was re-run successfully. A separate literal-rule implementation also
replayed all 1,957 distinct-symbol words and agreed on every one of the
19 legal-prefix states, the two successes, four stalls, and LCS length two.
This revision expands the literature audit and submission framing;
it does not assert a new mathematical theorem.
The original Prompt A and the full Prompt B input, audit and corrected
version-3 output are retained in
\texttt{research/pressing-3100089-pre-submission-review-2026-10-11/}.
The final publication audit, complete proposed and corrected source, and
repository-check outcomes are retained in
\texttt{research/pressing-3100089-submission-2026-10-11/}.
The same configured model performed these reviews. None constitutes
independent human verification. Only this version-3 manuscript and the
standalone verifier are included in the research PR; the earlier local
drafts and process transcripts are not publicly archived.

\section{Scope, remaining limits, and contributions}
\paragraph{Full-solution scope declaration.}
Theorem~\ref{counterexample} disproves the \emph{entire original universal
four-total-edit assertion} by one explicit admissible colored graph.
Every legal pressing branch on that graph is exhausted, and the entire
successful-word family has disconnected $H_4$. A counterexample is
sufficient to disprove the universal quantifier; it need not classify
all other graphs, colorings, empty families, isolated vertices, or the
empty graph. No significant external theorem or computation is a proof
dependency. There is no remaining mathematical gap in the disproof claimed
here. Independent human/formal verification and historical novelty are
unestablished. The distinct universal eight-total-edit formulation has
not been proved or disproved in this manuscript.

This manuscript therefore records \texttt{ATTEMPT\_STATUS: solved} as an
\emph{LLM full-disproof claim for the exact recorded target}. The catalogue
definition and its source status remain unchanged. Neither that metadata
nor the model's audit is an external acceptance or peer-review verdict.

\paragraph{Contributor declaration.}
The model-generated research was conducted under the guidance of
Wentao Li. OpenAI GPT 6.1 Sol, Max reasoning generated the search,
counterexample, proof, verification script, manuscript, and model-assisted
audits on 11 October 2026. Wentao Li set the research direction, supplied
iterative prompts, requested continued work toward a full resolution,
directed the proof and literature rechecks, and authorized this repository
submission. His contributor credit records research guidance, direction,
and prompting. No human mathematical proof development, substantial
mathematical correction, or independent human proof verification is recorded.
The earlier GPT-6-based Codex draft had its
exact version and reasoning unexposed; its attribution is preserved.
The double-star draft was by the same configured GPT 6.1 Sol, Max reasoning.
No other model or subagent materially contributed to this continuation.
Separate programs and repeated model review are model-assisted checks,
not independent human verification.

\section{Completion Estimate}
\noindent\textbf{COMPLETION: 100\%}
The original universal four-total-edit claim is fully disproved by
Theorem~\ref{counterexample}. The complete successful-word family of a
valid graph is disconnected, so no part of that universal assertion remains
to be settled for the claimed disproof. This is a scope-completion claim,
not proof confidence, historical novelty, independent verification, or
resolution of the distinct eight-total-edit formulation.

\begin{thebibliography}{99}
\bibitem{cooper} J. Cooper and J. Davis,
\emph{Successful Pressing Sequences for a Bicolored Graph and Binary Matrices},
arXiv:1502.07450v5, 2015, Definitions 1--2 and Conjecture 1, pp.~2--4.
\url{https://arxiv.org/abs/1502.07450v5}.
\bibitem{bixby} E. Bixby, T. Flint, and I. Mikl\'os,
\emph{Proving the pressing game conjecture on linear graphs},
arXiv:1303.6799, 2013, Conjecture 6 and Theorem 7;
Involve 9 (2016), 41--56.
\url{https://arxiv.org/abs/1303.6799}.
\url{https://doi.org/10.2140/involve.2016.9.41}.
\bibitem{jmm} B. Li, I. Mikl\'os, and C. R. Teplica,
\emph{The Connectedness of the Black-and-White Graph Pressing Game},
Joint Mathematics Meetings, San Francisco, 3--6 January 2024,
abstract 1192-05-32487, received 12 September 2023.
Indexed abstract consulted; full proof not accessed.
\url{https://www.ams.org/journals/abs/2024-45-01/1192.pdf}.
\bibitem{problems} J. Cooper, \emph{Combinatorial Problems I Like},
pressing-sequence connectivity entry, page last edited October 2020;
consulted 11 October 2026.
\url{https://people.math.sc.edu/cooper/combprob.html}.
\bibitem{unique} J. N. Cooper and H. W. Whitlatch,
\emph{Uniquely Pressable Graphs: Characterization, Enumeration, and Recognition},
arXiv:1706.07468. Accessible text consulted 11 October 2026.
\url{https://arxiv.org/abs/1706.07468}.
\bibitem{vposets} J. Cooper, P. Gartland, and H. Whitlatch,
\emph{A New Characterization of $\mathcal V$-Posets},
arXiv:1810.07276. Accessible text and Figure 2 inspected 11 October 2026.
\url{https://arxiv.org/abs/1810.07276}.
\bibitem{poster} J. Cooper and H. Whitlatch,
\emph{Genome Inversions and Graph Pressing Sequences}, research poster,
public author-hosted PDF, consulted 11 October 2026.
\url{https://connect.gonzaga.edu/asset/file/1363/Clemson_Poster.pdf}.
\bibitem{miklosthesis} I. Mikl\'os,
\emph{Computational complexity of counting and sampling problems in bioinformatics},
dissertation dc\_2046\_22, R\'enyi Institute, Chapter 9.
Conjecture statement and discussion inspected 11 October 2026.
\url{https://real-d.mtak.hu/1570/7/dc_2046_22_doktori_mu.pdf}.
\bibitem{whitlatchthesis} H. W. Whitlatch,
\emph{Successful Pressing Sequences in Simple Pseudo-Graphs},
Ph.D. dissertation, University of South Carolina, Spring 2019.
Landing page and abstract consulted; linked full text inaccessible.
\url{https://scholarcommons.sc.edu/etd/5135/}.
\bibitem{atlas} NetworkX developers, graph atlas data and documentation,
consulted 11 October 2026; based on R. C. Read and R. J. Wilson,
\emph{An Atlas of Graphs}, Oxford University Press, 1998.
Discovery data only; not a dependency of the disproof.
\url{https://networkx.org/documentation/stable/reference/generated/networkx.generators.atlas.graph_atlas_g.html}.
\url{https://github.com/networkx/networkx/blob/main/networkx/generators/atlas.dat.gz}.
\bibitem{priorstar} OpenAI GPT-6-based Codex, exact version and reasoning
unexposed, \emph{Successful pressing sequences on stars and star forests},
11 October 2026, preserved local unsubmitted draft:
\texttt{attacks/open\_problems/top\_problems/}
\texttt{gpt6\_codex\_version\_unexposed/3100089.tex}.
\bibitem{priordouble} OpenAI GPT 6.1 Sol, Max reasoning,
\emph{A finite reduction and certificate for four-edit connectivity on double stars},
11 October 2026, preserved local unsubmitted draft:
\texttt{attacks/open\_problems/top\_problems/gpt6\_1\_sol\_max/3100089.tex}.
\bibitem{catalogue} ULAM.AI and the UnsolvedMath Contributors,
\emph{UnsolvedMath}, record 3100089 (AMR-030-0089), and the repository
definition at commit b1f633121074e16579fa7fd05314e5f504d29bae.
Original curation and metadata: CC BY 4.0.
\url{https://huggingface.co/datasets/ulamai/UnsolvedMath}.
\url{https://github.com/mehmetmars7/Erdosproblems-llm-hunter/blob/b1f633121074e16579fa7fd05314e5f504d29bae/attacks/open_problems/top_problems/definitions/3100089.tex}.
\url{https://creativecommons.org/licenses/by/4.0/}.
\end{thebibliography}
\end{document}
